\documentclass[11pt]{article}
\usepackage[utf8]{inputenc}
\usepackage{amsmath,amssymb,amsfonts}
\usepackage{graphicx}
\usepackage{booktabs}
\usepackage{algorithm}
\usepackage{algorithmic}
\usepackage{natbib}
\usepackage[margin=1in]{geometry}
\usepackage{hyperref}

% Utility commands
\newcommand{\ie}{i.e.}
\newcommand{\eg}{e.g.}
\newcommand{\cf}{cf.}

\title{Thermodynamic Memory vs.\ Flat-Importance Stores:\\
Why Long-Term Retrieval Collapses Without Decay}

\author{
  Clement Deust\\
  Independent Researcher\\
  \texttt{github.com/cdeust/Cortex}
}

\date{April 2026}

\begin{document}
\maketitle

\begin{abstract}
External memory for large language models is dominated by
\emph{flat-importance} stores: vector indexes, BM25 corpora, and
long-context buffers in which every item carries the same long-term
retrieval prior.  We argue this design is asymptotically broken.
As the corpus grows, top-$k$ retrieval over an undifferentiated pool
degenerates into near-arbitrary tie-breaking among items with
comparable surface similarity, and the discriminative information
delivered to the consumer LLM approaches zero---a failure mode that
compounds the position bias of long-context decoding
\citep{Liu2023}.  We formalise the flat-importance failure, then
describe Cortex, a memory architecture that maintains a non-flat
priority distribution across $N$ by coupling four mechanisms:
(i)~continuously decaying \emph{heat} on every item
\citep{Ebbinghaus1885}, (ii)~a hierarchical predictive-coding write
gate \citep{Friston2010}, (iii)~consolidation cascades that compress
episodic into semantic memory \citep{Kandel2001,McClelland1995}, and
(iv)~WRRF fusion with heat as a tie-breaker.  On three independent
long-term-memory benchmarks, Cortex reaches
LongMemEval R@10\,$=$\,97.8\% (vs.\ 78.4\% paper-best),
LoCoMo R@10\,$=$\,92.6\%, and BEAM Overall\,$=$\,0.543
(vs.\ 0.329 paper-best).  We discuss when flat memory remains adequate
(small $N$, single-session contexts), the calibration cost of decay,
and the per-write overhead of biological consolidation.  Code, data,
and benchmark harness are publicly available.\footnote{\url{https://github.com/cdeust/Cortex}}
\end{abstract}

%----------------------------------------------------------------------
\section{Introduction}
\label{sec:intro}

The default external memory for LLM agents is some variant of the
same recipe: embed the input, write to a vector index, and at query
time retrieve the top-$k$ nearest neighbours, optionally re-ranked.
This recipe has succeeded at the document-retrieval scale it was
designed for \citep{Lewis2020}.  It does not scale to the use case
agentic systems are pushing it toward---open-ended, multi-month,
mixed-content memory across thousands of conversations.

The reason is not retrieval speed but \emph{retrieval semantics}.
A vector store treats every entry as equally important in the long
term.  There is no mechanism by which a fact mentioned in passing
six months ago is given less retrieval weight than a decision the
user explicitly anchored last week.  As $N$ grows, the top-$k$ for a
typical query fills with semantically plausible but contextually
irrelevant items---and the downstream LLM, subject to the position
bias documented in \emph{Lost in the Middle} \citep{Liu2023},
processes them in an order uncorrelated with their actual usefulness.

This paper argues that the missing mechanism is decay, and that
decay is not a heuristic but a structural requirement: it is what
keeps the priority distribution over memories non-flat at any $N$.
We characterise the flat-importance failure mode, describe Cortex's
thermodynamic architecture, and report empirical results on three
published long-term-memory benchmarks.

\paragraph{Teaser.}
On LongMemEval \citep{Wu2025}, Cortex reaches Recall@10 of 97.8\%,
vs.\ the paper-best of 78.4\%.  On BEAM \citep{Tavakoli2026}, Cortex
reaches 0.543 Overall, vs.\ 0.329 paper-best.  All numbers are from
runs against Cortex's production database, single-process, with the
exact PL/pgSQL retrieval code path used in deployment.

%----------------------------------------------------------------------
\section{Background and Related Work}
\label{sec:related}

\paragraph{Retrieval-augmented generation.}
The canonical RAG architecture \citep{Lewis2020} couples a dense
retriever to a generator.  Document encoders such as DPR
\citep{Karpukhin2020} and contrastive bi-encoders dominate the
retrieval side.  These systems target a fixed, curated corpus where
every document is, by construction, equally part of the index.
Long-term memory is not part of their formulation.

\paragraph{Long-context LLMs.}
A parallel line of work scales the context window
itself---200k tokens (Anthropic Claude 2.1, 2023), 1M+ tokens
(Google Gemini 1.5, 2024).  A long context substitutes for retrieval
by holding everything in attention.  This works until the relevant
fact is buried in the middle of the window, at which point
\citet{Liu2023} show retrieval accuracy degrades sharply: U-shaped
position bias persists across model families and scales.  Long
context does not solve discriminability---it inherits it.

\paragraph{Episodic memory in LLM agents.}
MemGPT \citep{Packer2024} introduced a hierarchical paged memory
metaphor with explicit eviction.  Generative Agents \citep{Park2023}
added importance scores from a separate model call and reflection
passes that summarised episodes into higher-level conclusions.  Both
are first steps toward non-flat memory: MemGPT decides what stays
in the active page; Generative Agents weight retrieval by importance
$\times$ recency $\times$ similarity.  Neither implements the full
biological cascade---predictive-coding gating, multi-level
consolidation, neuromodulatory tagging, microglial pruning---that
we draw on here.

\paragraph{Long-term-memory benchmarks.}
LongMemEval \citep{Wu2025}, LoCoMo \citep{Maharana2024}, and BEAM
\citep{Tavakoli2026} each isolate the long-horizon retrieval failure
mode that \citet{Liu2023} anticipated.  We report on all three in
\S\ref{sec:empirical}.

\paragraph{Cortex ecosystem.}
This paper describes the memory engine of a four-component
ecosystem: \emph{cortex-beam-abstain} \citep{Deust2026a} is a
learned retrieval-abstention model that complements decay on
residual collapse; \emph{automatised-pipeline} \citep{Deust2026b} is
a Rust tree-sitter / LadybugDB property-graph backend that anchors
memories to code symbols; \emph{prd-spec-generator}
\citep{Deust2026c} is the read-heavy downstream consumer that
justifies the per-write cost; \emph{zetetic-team-subagents}
\citep{Deust2026d} is the 116-agent reasoning library that drives
the workload distribution.  \S\ref{sec:ecosystem} expands on each.

%----------------------------------------------------------------------
\section{The Flat-Importance Failure Mode}
\label{sec:flat-failure}

\subsection{Definition}
\label{sec:flat-def}

A memory store is \emph{flat-importance} if, for any two stored
items $x_i, x_j$ at time $t$, the long-term retrieval prior is
independent of access history, recency, surprise, or causal context:
\begin{equation}
P(x_i \mid t) = P(x_j \mid t) \quad \text{for all } i, j \text{ once written.}
\label{eq:flat-prior}
\end{equation}
A canonical vector store is flat-importance: every entry contributes
a single embedding, and only query-conditional similarity differs
at retrieval time.

\subsection{Asymptotic Collapse}
\label{sec:flat-collapse}

Under flat priors, the retrieval score for a query $q$ depends only
on similarity $s(q, x_i)$.  Let $S_q$ be the empirical distribution
of $s(q, x_i)$ over the corpus.  As $N$ grows, two effects compound:
\begin{enumerate}
  \item \textbf{Concentration.}  With high-dimensional embeddings,
        the gap between the top-1 and top-$k$ similarity scores
        shrinks \citep{Aggarwal2001,Beyer1999}.  The top-$k$ pool
        fills with near-tied items.
  \item \textbf{Tie-breaking is uninformed.}  Without a
        non-similarity prior, ordering among near-tied items is
        determined by floating-point noise and storage-order
        artefacts.
\end{enumerate}

A useful informal statement:
\begin{quote}
\emph{Claim (flat collapse).}  If the retrieval prior $P(x)$ is
constant and $s(q, x)$ is the only ranking signal, then as
$N \to \infty$ the conditional entropy
$H(\mathrm{useful} \mid \mathrm{retrieved\ top}\text{-}k) \to H(\mathrm{useful})$;
that is, the retrieval step delivers no information about which
retrieved items are actually useful for $q$.
\end{quote}

This is the \emph{epistemic} failure mode.  It is independent of
and prior to any LLM consumption issue.  A high-dimensional
hubness analysis of the same regime is given by
\citet{Radovanovic2010}.

\subsection{Consequences for the LLM Consumer}
\label{sec:flat-consumer}

Even if retrieval merely returned an unordered set, the position
bias documented by \citet{Liu2023} would still degrade downstream
answers when $k$ is large.  With retrieval order itself
uninformative, the position bias compounds the flat-prior problem:
a fact that \emph{is} useful gets sub-ranked into the U-shaped
middle and is effectively ignored.

This is why na\"{i}ve ``stuff more into context'' approaches plateau,
and why ``raise $k$'' is not a fix at large $N$: the information
density of retrieved tokens has dropped in proportion.

%----------------------------------------------------------------------
\section{Cortex's Thermodynamic Memory Model}
\label{sec:thermodynamic}

Cortex's architecture is an attempt to keep the retrieval prior
$P(x)$ non-flat at every $N$ by implementing the mechanisms
biological memory uses for the same problem.  Each mechanism cites
a primary source and a Cortex implementation file.

\subsection{Heat as a Continuous Priority Field}
\label{sec:heat}

Every memory carries a scalar \texttt{heat} in $[0, 1]$.  Heat
decays exponentially over time and is reheated on access---the
\citet{Ebbinghaus1885} forgetting curve, $R(t) = e^{-t/S}$, where
$S$ is memory stability.

Cortex's decay
(\texttt{mcp\_server/core/thermodynamics.py::compute\_decay}) is
\begin{equation}
h(t) = h_0 \cdot \lambda_{\mathrm{eff}}^{\,t}
\label{eq:heat-decay}
\end{equation}
with the per-hour decay factor $\lambda_{\mathrm{eff}}$ modulated by
importance and valence:
\begin{equation}
\lambda_{\mathrm{eff}} = 1 - \frac{1 - \lambda_{\mathrm{base}}}
                                     {1 + |v| \cdot r \cdot (1 - e^{-t})}
\label{eq:lambda-eff}
\end{equation}
where $\lambda_{\mathrm{base}} = 0.95$ for normal memories,
$0.998$ for high-importance memories ($I > 0.7$), $r = 0.5$ is the
emotional resistance, and the time-saturation factor $1 - e^{-t}$
implements the \citet{Kleinsmith1963} crossover effect: emotional
advantage grows with delay \citep{Yonelinas2015}.  Half-lives:
$\approx 14$\,h for normal items, $\approx 346$\,h for important
items.  The constants are calibrated to a hours/days timescale and
are not from a single paper; they are documented as engineering
choices in the source.

\subsection{Predictive-Coding Write Gate}
\label{sec:gate}

Not every input is written.  Cortex routes incoming content
through a hierarchical \citet{Friston2010} free-energy gate
(\texttt{mcp\_server/core/hierarchical\_predictive\_coding.py})
operating at three levels---sensory, entity, schema---analogous to
the cortical hierarchy.  Items that fail to surprise the model at
any level are discarded; items that surprise at the schema level
receive a higher initial heat.  This implements a write-side
scarcity: the store grows by what is actually informative, not by
what merely passes through.

\subsection{Coupled Neuromodulation}
\label{sec:neuromod}

Heat is not the only modulator.  Cortex implements a coupled
DA / NE / ACh / 5-HT cascade
(\texttt{mcp\_server/core/coupled\_neuromodulation.py}) with
cross-channel effects, following \citet{Doya2002} and
\citet{Schultz1997}.  Dopamine prediction errors raise importance
for items whose outcome surprised the prior; norepinephrine
reshapes the gain on uncertainty.

\subsection{Emotional Tagging}
\label{sec:emotional}

\citet{Yerkes1908} characterised the inverted-U relationship
between arousal and performance.  \citet{Wang2024} describe
amygdala-priority encoding for emotionally salient inputs.
Cortex's \texttt{mcp\_server/core/emotional\_tagging.py} combines
both: arousal modulates the priority of incoming memories, and the
inverted-U bound prevents over-weighting.

\subsection{Synaptic Tagging}
\label{sec:synaptic}

\citet{Frey1997} showed that weak memories sharing entities with
later strong potentiation events can be retroactively promoted.
Cortex implements the analogue in
\texttt{mcp\_server/core/synaptic\_tagging.py}: when a
strongly-encoded memory arrives, weak memories sharing entities
receive a heat boost, preserving items the system did not yet know
were important.

\subsection{Consolidation Pipeline}
\label{sec:consolidation}

Episodic memories progress through a four-stage cascade---LABILE
$\to$ EARLY\_LTP $\to$ LATE\_LTP $\to$ CONSOLIDATED---modelled on
\citet{Kandel2001}.  Implementation:
\texttt{mcp\_server/core/cascade.py}.  Once consolidated, the
Complementary Learning Systems theory of \citet{McClelland1995}
governs the episodic$\to$semantic transfer in
\texttt{mcp\_server/core/dual\_store\_cls.py}: hippocampal-style
episodic traces feed cortical-style schemas, and the episodic
surface area shrinks while structured knowledge accumulates.  See
also \citet{Tonegawa2015} on the cellular substrate (engram cells)
and \citet{Hopfield1982} on the content-addressable view of
attractor memory.

\subsection{Sleep Replay}
\label{sec:sleep}

\citet{Buzsaki2015} and \citet{Hasselmo2005} describe phase-gated
encoding/retrieval/consolidation cycles.  Cortex's
\texttt{oscillatory\_clock.py} and \texttt{replay.py} implement an
offline replay pass that consolidates and re-embeds clusters during
idle windows.

\subsection{Pattern Separation and Pruning}
\label{sec:separation}

\citet{Leutgeb2007} and \citet{Yassa2011} describe dentate-gyrus
pattern separation.  Cortex's \texttt{pattern\_separation.py}
orthogonalises near-duplicate items so they remain distinguishable.
\citet{Wang2020} describe complement-dependent microglial pruning;
Cortex's \texttt{microglial\_pruning.py} prunes the lowest-utility
edges of the entity graph, reclaiming structure rather than mass.

%----------------------------------------------------------------------
\section{Why Decay Prevents Collapse}
\label{sec:why-decay}

The load-bearing argument of this paper is that the mechanisms in
\S\ref{sec:thermodynamic} keep the retrieval prior $P(x)$ non-flat
at all $N$, and that this is what preserves discriminability.  We
make this argument in three forms.

\subsection{Information-Theoretic}
\label{sec:info-theory}

If heat decays exponentially and is reheated on access, then in
steady state the heat distribution across an active corpus is
well-approximated by a power law (long tail of cold items, short
head of repeatedly-accessed hot items).  Define the
priority-weighted retrieval entropy
\begin{equation}
H_P(R \mid q) \;=\; -\sum_{x \in R}
\frac{P(x)\, s(q, x)}{Z}
\log \frac{P(x)\, s(q, x)}{Z},
\quad Z = \sum_{x \in R} P(x)\, s(q, x).
\label{eq:HP}
\end{equation}
Under flat $P(x) = 1/N$, $H_P$ collapses onto the similarity-only
distribution and inherits its high-dimensional concentration
(\S\ref{sec:flat-collapse}).  Under a power-law $P(x)$, $H_P$
retains positive mass on the head even when the similarity
distribution is near-tied: the tie-breaker is the heat prior itself.
We do not claim this as a tight bound; we claim it as the
\emph{mechanism} by which decay preserves
$H(\mathrm{useful} \mid \mathrm{retrieved})$ above the flat-prior
floor.

\subsection{Operational---WRRF Fusion}
\label{sec:wrrf}

The Cortex retrieval path is implemented server-side in PL/pgSQL
(\texttt{recall\_memories()}, see
\texttt{mcp\_server/infrastructure/pg\_schema.py}).  It computes a
\emph{Weighted Reciprocal Rank Fusion} \citep{Cormack2009} over six
independent signals:
\begin{equation}
\mathrm{score}(x \mid q) \;=\;
\sum_{c \in \{\mathrm{vec}, \mathrm{fts}, \mathrm{tri}, \mathrm{heat}, \mathrm{rec}, \mathrm{ngram}\}}
\frac{w_c}{k + \mathrm{rank}_c(x \mid q)}
\label{eq:wrrf}
\end{equation}
with intent-conditioned weights $w_c$ chosen by the query router
(\texttt{mcp\_server/core/query\_router.py}).  When two items have
equivalent vector and FTS scores, the heat and recency channels
resolve the tie---so the priority field $P(x)$ acts not just as a
prior but as a deterministic disambiguator in the regime where
similarity is uninformative.  This is the operational mechanism
that prevents \S\ref{sec:flat-collapse}'s tie-noise.

\subsection{Curation by Attrition}
\label{sec:curation}

Decay also removes items, but only after they have been compressed.
The CLS pipeline (\texttt{dual\_store\_cls.py}) detects clusters of
cold-but-related episodic items, abstracts them into semantic
schemas, and lets the underlying episodic surface fade.  Two
consequences: (a)~the retrievable corpus stays small relative to
the input stream, so high-dimensional concentration is held off;
(b)~the knowledge that mattered survives in compressed form even
when the verbatim episodes do not.

The combination---write-side scarcity (\S\ref{sec:gate}),
heat-modulated tie-breaking (\S\ref{sec:wrrf}), and
curation-by-attrition (\S\ref{sec:curation})---is what we mean by
``non-flat at any $N$.''  None of the three on its own would
suffice.

\subsection{Decay Reduces the \emph{Rate} of Collapse, Not Its Possibility}
\label{sec:rate-not-possibility}

The mechanisms above reduce how often retrieval enters a regime
where similarity scores are uninformative; they cannot guarantee the
answer is in the store.  For queries whose target is genuinely
absent, the WRRF tie-breaker simply chooses the most-prior of a set
of irrelevant candidates---the right response is abstention, not
retrieval.

The diagnostic data published alongside cortex-beam-abstain
\citep{Deust2026a}, corroborated by BEAM \citep{Tavakoli2026},
makes the point sharply: on BEAM, abstention queries (questions
whose answer is not in the haystack) receive an \emph{average
retrieval similarity score of 0.926}---higher than many answerable
queries.  Score-based thresholding cannot decide abstention on this
distribution; the relevance score is doing the wrong job.  Cortex
reduces the \emph{frequency} of discriminability collapse via decay;
cortex-beam-abstain (a learned binary classifier;
\S\ref{sec:ecosystem}) handles the \emph{residual} collapse via
learned abstention.  The two mechanisms are complementary: the
0.926 figure is direct empirical evidence that flat-importance
$+$ score-based-threshold retrieval is fundamentally insufficient.

%----------------------------------------------------------------------
\section{Empirical Evidence}
\label{sec:empirical}

We evaluate Cortex on three independent long-term-memory benchmarks.
All numbers are from runs against the production PostgreSQL database
with the same \texttt{recall\_memories()} PL/pgSQL code path used in
deployment, single-process, on a clean DB seeded only with each
benchmark's data.

\begin{table}[t]
\centering
\caption{Cortex vs.\ paper-best on three long-term-memory benchmarks.
Cortex numbers: clean DB, single process, April 2026.}
\label{tab:benchmarks}
\begin{tabular}{llrrr}
\toprule
Benchmark & Venue & Metric & Cortex & Paper-best \\
\midrule
LongMemEval & ICLR 2025 & R@10    & \textbf{97.8\%} & 78.4\% \\
LongMemEval & ICLR 2025 & MRR     & \textbf{0.882}  & ---    \\
LoCoMo      & ACL 2024  & R@10    & \textbf{92.6\%} & ---    \\
LoCoMo      & ACL 2024  & MRR     & \textbf{0.794}  & ---    \\
BEAM        & ICLR 2026 & Overall & \textbf{0.543}  & 0.329  \\
\bottomrule
\end{tabular}
\end{table}

The $+19.4$\,pp absolute gain on LongMemEval R@10 and the $+65\%$
relative gain on BEAM Overall are the headline results.  Both
benchmarks include question categories specifically designed to
defeat flat retrieval---multi-session reasoning (LongMemEval),
causal/temporal grounding (BEAM)---which is consistent with the
\S\ref{sec:flat-collapse} claim that the flat regime fails fastest
on questions that require integrating information across the
priority distribution.

\paragraph{Where Cortex wins.}
The largest gaps appear on temporal questions (``what did I decide
first about X?''), causal-chain questions (``why did Y change?''),
and multi-hop knowledge integration.  These are the question
categories that require traversal of the entity/causal graph and
that benefit most from heat-modulated tie-breaking.

\paragraph{Where the gap is smaller.}
On surface-fact retrieval (``what is the value of X?'') with a
small corpus, flat baselines do reasonably well---there is no
priority disambiguation to do because the fact is uniquely
identified by similarity.  Cortex's advantage on these categories
is primarily from the FTS/trigram channels, not the thermodynamic
ones.

\paragraph{Caveats on these numbers.}
(i)~We do not have head-to-head re-runs of every published baseline
on our exact protocol; we report Cortex's numbers and the highest
paper-reported number on each benchmark.
(ii)~These benchmarks are retrieval-quality benchmarks; downstream
end-task accuracy with a specific LLM may differ.
(iii)~BEAM's Overall is a composite of seven sub-metrics---see
\texttt{benchmarks/beam/} for the per-subset breakdown.

%----------------------------------------------------------------------
\section{Discussion}
\label{sec:discussion}

\subsection{Limitations}
\label{sec:limitations}

\paragraph{Decay calibration.}
The constants in \S\ref{sec:heat} ($\lambda_{\mathrm{base}} = 0.95$,
$\lambda_{\mathrm{important}} = 0.998$, $r = 0.5$) are calibrated
to a hours/days timescale and are explicitly engineering choices,
not paper-derived.  Too aggressive a forgetting curve loses
cold-but-useful facts; too gentle a curve recovers the flat-prior
failure mode.  Cortex pins the calibration with \texttt{rate\_memory}
user feedback, which feeds metamemory confidence
(\texttt{compute\_metamemory\_confidence}), but this is a
per-deployment calibration loop, not a universal solution.

\paragraph{Cold-start regime.}
The argument in \S\ref{sec:flat-collapse} is asymptotic.  For small
corpora ($N < {\sim}10$k items), high-dimensional concentration has
not yet bitten, and a well-tuned flat RAG is competitive.  Cortex's
advantage grows with $N$; we do not claim it dominates at every
scale.

\paragraph{Per-write cost.}
The thermodynamic write path is heavier than appending to a vector
index: predictive-coding gate, emotional tagging, synaptic tagging
propagation, cascade-stage assignment, entity extraction.  In our
deployment writes are roughly $100\times$ rarer than reads, so the
amortised cost is acceptable; in write-heavy workloads it would not
be.

\paragraph{Coverage.}
We have not run Cortex against MemoryAgentBench or EverMemBench
yet; those are next on the roadmap (\texttt{benchmarks/memoryagentbench/},
\texttt{benchmarks/evermembench/}).

\paragraph{Ablation completeness.}
Cortex includes an ablation framework (\texttt{mcp\_server/core/ablation.py})
for 23 mechanisms, but the per-mechanism contribution to the
headline scores has not been fully reported in this paper.  A full
ablation (heat-only, gate-only, CLS-only, full system) would
tighten the causal story for \S\ref{sec:why-decay}; currently we
report only the integrated system.

\subsection{When Flat Memory is Fine}
\label{sec:when-flat}

\begin{itemize}
  \item Single-session, single-document tasks (the original RAG
        setting).
  \item Curated corpora that are themselves the result of a
        non-flat curation process upstream.
  \item Workloads where $N$ is bounded by construction (a fixed
        knowledge base).
  \item Latency-critical paths where the per-write cost of the
        predictive-coding gate is unacceptable.
\end{itemize}

The argument of this paper is not that flat is wrong, but that it
is wrong \emph{as the long-term memory of an agent that lives
indefinitely}.

\subsection{What This Paper is Not}
\label{sec:not}

\begin{itemize}
  \item We do not claim Cortex ``solves'' memory.  The above
        limitations are real.
  \item We do not compare against baselines whose numbers we do
        not have.  The ``paper-best'' column is the highest
        published number we found on each benchmark.
  \item We do not claim biological fidelity; we claim biological
        \emph{inspiration}, with each module documenting which
        paper it draws on and where it deviates.
\end{itemize}

%----------------------------------------------------------------------
\section{The Cortex Ecosystem: Where This Memory Model Fits}
\label{sec:ecosystem}

The thermodynamic memory model described above is calibrated to a
specific workload---agentic coding sessions issuing many recall
calls per write against memories that link to code structure---and
to a deployment shape in which four other software components
handle cases this paper's mechanisms are not designed to solve.

\paragraph{cortex-beam-abstain---the abstention complement.}
\S\ref{sec:rate-not-possibility} establishes that decay reduces the
\emph{rate} of discriminability collapse but cannot eliminate it
when the answer is absent.  cortex-beam-abstain \citep{Deust2026a}
is a DistilBERT binary classifier (66M params, INT8 ONNX, 64\,MB,
${\sim}32$\,ms per pair on CPU) trained on 19{,}111 (query, passage)
pairs from BEAM splits 100K/500K/1M with hard negatives mined by
query-passage cosine, reaching $F_1 = 0.733$ on the v0.1 release
(score range $0.215$--$0.830$).  It is the \emph{complement} to
decay: Cortex addresses the \emph{frequency} of collapse,
cortex-beam-abstain addresses the \emph{residual}.  Both mechanisms
are required; neither suffices alone.

\paragraph{automatised-pipeline---the code-structural anchor.}
Cortex memories are not anonymous text snippets; downstream
consumers expect them to refer to specific functions, types, and
call paths.  automatised-pipeline \citep{Deust2026b} is a Rust MCP
server that indexes Rust / Python / TypeScript codebases via
tree-sitter into a LadybugDB property graph (16 node labels, 36+
relationship tables; 220 tests passing) and exposes 23 MCP tools
for symbol resolution, impact estimation, and PRD-vs-graph
validation.  This lets Cortex's memories link to actual symbols
rather than substrings---recall returns text-grounded-in-structure,
not text alone.

\paragraph{prd-spec-generator---the read-heavy downstream consumer.}
\S\ref{sec:limitations} noted that writes in our deployment are
roughly $100\times$ rarer than reads.  prd-spec-generator
\citep{Deust2026c} is the canonical example: a TypeScript MCP
server (10 packages, 17 MCP tools, 583 tests) that turns a feature
description into a 9-file PRD via a stateless reducer.  Each
section synthesis issues multiple Cortex recall calls; a typical
trial-tier feature run produces ${\sim}62$ host-visible iterations,
each with its own retrieval.  This is the workload that amortises
Cortex's thermodynamic write cost.

\paragraph{zetetic-team-subagents---the workload distribution.}
The 116-agent reasoning library (97 genius $+$ 19 team agents;
241 memory tests passing) drives the queries Cortex sees in
production \citep{Deust2026d}.  Multi-agent chains---\texttt{peirce
$\to$ cochrane $\to$ feynman $\to$ toulmin} for deep research,
\texttt{fermi $\to$ curie $\to$ knuth} for performance
investigation---generate clusters of related recalls within a
session, shaping the query distribution against which Cortex's
decay constants and intent profiles are calibrated.  The same
library imposes the source-citation discipline this paper itself
follows.

These are not afterthoughts; they shape what the memory store has
to be good at.  A memory system optimised for a different workload
---interactive conversation logs without code or PRDs,
single-session tutoring with no read-ahead reuse---would calibrate
the decay rate, the predictive-coding gate, and the consolidation
pipeline differently.  \textbf{The thermodynamic model is general;
its constants are not.}  A SQLite variant
(Cortex-cowork, \citealp{Deust2026e}) ports the same architecture
to sandboxed environments lacking PostgreSQL---client-side fusion
replaces PL/pgSQL, \texttt{sqlite-vec} flat scan replaces pgvector
HNSW---but the decay model and consolidation pipeline are
unchanged, consistent with the claim that the model is
implementation-independent.

%----------------------------------------------------------------------
\section{Conclusion}
\label{sec:conclusion}

Flat-importance memory stores are not a stable long-term substrate
for agentic LLMs.  As $N$ grows, the absence of a non-similarity
prior collapses top-$k$ retrieval into noise tie-breaking, and the
position bias of long-context decoding compounds the loss.  Cortex
addresses this by maintaining a non-flat priority distribution over
memories---implemented through heat decay, a hierarchical
predictive-coding write gate, neuromodulatory tagging, and a
four-stage consolidation cascade---and by using that distribution
as both a retrieval prior and a tie-breaker in WRRF fusion.  On
three long-term-memory benchmarks, Cortex outperforms the published
bests by margins consistent with the theoretical argument: largest
on the question categories that most require non-flat priors,
smaller on surface-fact retrieval where similarity alone suffices.

The ablation work needed to make this story tight---per-mechanism
contribution to each benchmark, dose--response curves for the decay
constants, head-to-head re-runs of published baselines on the same
protocol---is the natural next step.

%----------------------------------------------------------------------
\appendix

\section{Information-Theoretic Collapse of Flat-Importance Retrieval}
\label{app:shannon}

This appendix derives why a \emph{flat-importance} memory store
loses retrieval discriminability as $N$ grows, while a
\emph{decaying-priority} store (Cortex) preserves it.  The argument
follows \citet{Shannon1948}: define the right quantity, derive its
limit, then compare designs.

\subsection{Setup and Notation}
\label{app:shannon-setup}

A memory store holds $N$ items $\{x_1, \dots, x_N\}$.  Each item
$x_i$ carries an \emph{importance} $w_i \ge 0$ with $\sum_i w_i = 1$.
Given a query $q$, the retrieval score is
\begin{equation}
s_i(q) \;=\; \alpha \cdot \mathrm{sim}(q, x_i) \;+\;
             (1 - \alpha) \cdot w_i, \qquad \alpha \in [0,1].
\label{eq:score-mix}
\end{equation}
Top-$k$ retrieval returns the indices
$R_k(q) = \arg\text{top-}k_i\, s_i(q)$.

\textbf{Two regimes.}
\begin{itemize}
  \item \textbf{Flat-importance (canonical RAG).}
        $w_i = 1/N$ for all $i$.  The importance term is constant,
        so ranking reduces to ranking $\mathrm{sim}(q, x_i)$.
  \item \textbf{Decaying-priority (Cortex).}
        $w_i \propto h_i \cdot \exp(-\lambda \,(t - t_i))$, where
        $h_i$ is a base heat (access frequency) and $t - t_i$ is
        the age since last set.  The exponential form is
        Ebbinghaus's forgetting law \citep{Ebbinghaus1885},
        recoverable as the steady state of a Friston-style
        free-energy minimization with linear precision decay
        \citep{Friston2010}.
\end{itemize}

We treat $\mathrm{sim}(q, x_i)$ as random variables drawn iid from a
distribution $F_S$ with mean $\mu_S$, variance $\sigma_S^2 < \infty$,
and smooth density $f_S$ on the relevant support.  This is the
iid-sim assumption (revisited as a limitation in
\S\ref{app:shannon-limit}).

\subsection{Discriminability Metric}
\label{app:shannon-disc}

Define the \emph{top-$k$ discriminability} as the mutual information
between the retrieved set $R_k$ and the query $q$:
\begin{equation}
D_k \;=\; I(R_k; Q) \;=\; H(R_k) - H(R_k \mid Q).
\label{eq:Dk}
\end{equation}
When $D_k \to 0$, the retrieved set is independent of the
query---the system has \emph{collapsed} into noise.  $D_k$ is
upper-bounded by $\log \binom{N}{k}$ and is the operational
quantity for ``does this retriever discriminate?'' in the Shannon
sense: the limit of bits-per-query the channel from query to
retrieved set can carry.

\subsection{The Collapse Theorem (Informal)}
\label{app:shannon-collapse}

\textbf{Claim.} Under flat importance ($w_i = 1/N$), iid
similarities with finite variance $\sigma_S^2$ and smooth density
$f_S$, and embedding noise floor $\eta$, the expected gap between
the $k$-th and $(k+1)$-th largest score satisfies
\[
\mathbb{E}\big[s_{(k)} - s_{(k+1)}\big]
\;\xrightarrow[N \to \infty]{}\; 0,
\]
and the top-$k$ ranking becomes effectively random once the gap
drops below $\eta$.  Consequently $D_k \to 0$.

\paragraph{Derivation sketch (order statistics).}
Let $S_{(1)} \ge S_{(2)} \ge \dots \ge S_{(N)}$ denote the ordered
similarities.  From \citet{David2003}, \S4.6, for a smooth density
$f_S$ and a quantile $p_k = 1 - k/N$ with $k$ fixed and $N$ large:
\begin{equation}
\mathbb{E}\big[S_{(k)} - S_{(k+1)}\big]
\;\approx\;
\frac{1}{N \, f_S\!\big(F_S^{-1}(p_k)\big)}
\;=\; O\!\big(1/N\big),
\label{eq:order-spacing}
\end{equation}
with standard deviation of $S_{(k)}$ itself scaling as
$\sigma_S / \sqrt{N}$ around its quantile.  The \emph{spacing}
(which determines tie-breaking) shrinks faster than the position
uncertainty---both vanish, but the spacing collapses first.

\paragraph{Embedding noise floor.}
Cosine similarity on high-dimensional sentence-transformer
embeddings has empirical reproducibility $\eta \approx 0.05$
\citep{Reimers2019}.  Once
\[
\frac{1}{N \, f_S(F_S^{-1}(p_k))} \;<\; \eta,
\]
the rank order of $S_{(k)}, S_{(k+1)}$ is dominated by embedding
noise rather than semantic signal.  Tie-breaking becomes uniform,
and $H(R_k \mid Q) \to H(R_k)$, hence $D_k \to 0$.

\paragraph{Threshold.}
Setting the spacing equal to $\eta$ gives the critical store size
\begin{equation}
N^\star \;=\; \frac{1}{\eta \, f_S(F_S^{-1}(p_k))}.
\label{eq:Nstar}
\end{equation}
For $\eta = 0.05$ and a typical bulk density $f_S \approx 4$
(sentence-transformer cosines concentrate in $[0.2, 0.6]$),
$N^\star \approx 5{,}000$.  Beyond that, flat-importance retrieval
enters the collapse regime.

\subsection{Why Decay Restores Discriminability}
\label{app:shannon-restore}

\textbf{Power-law importance.} Empirical access patterns in human
and artificial memory follow Zipf-like laws: \citet{Anderson1989}
and \citet{Anderson1991} showed that the probability a memory is
needed at time $t$ scales as
$P(\mathrm{needed}) \propto t^{-\gamma}$ with
$\gamma \in [0.5, 1.5]$ across email, headlines, and child-directed
speech corpora.  The power law of forgetting \citep{Wixted1991}
gives the same shape.  When Cortex's heat
$h_i \cdot \exp(-\lambda (t - t_i))$ is averaged over the access
process, the stationary distribution of $w_i$ is heavy-tailed with
tail exponent $\gamma$.

\textbf{Score distribution becomes bimodal.}  The composite score
$s_i = \alpha \cdot \mathrm{sim}(q, x_i) + (1 - \alpha) \cdot w_i$
is a convolution of a thin-tailed similarity component with a
heavy-tailed importance component.  For $\alpha \in [0.3, 0.7]$
(Cortex's WRRF operating regime), the resulting distribution is
bimodal: a sharp head from items with large $w_i$ and a long tail
of low-importance items.  Top-$k$ retrieval consistently returns
from the head; tie-breaking is dominated by $w_i$, which has gaps
$O(1)$---not by sub-$\eta$ similarity differences.

\textbf{Non-vanishing entropy bound.}  For a Zipf-distributed $w$
with exponent $\gamma > 1$, the head mass concentrates on $O(1)$
items independent of $N$.  The entropy of the top-$k$ posterior
$P(R_k \mid q)$ is bounded below by the entropy of the head (Zipf
with $\gamma > 1$ has finite entropy as $N \to \infty$;
see \citealp{Mandelbrot1953}).  Therefore
\begin{equation}
\liminf_{N \to \infty} D_k \;\ge\; H_{\mathrm{head}}(\gamma, k)
\;>\; 0.
\label{eq:Dk-bound}
\end{equation}
Decay is what generates and \emph{maintains} the heavy tail:
without it, repeated writes regress $w_i$ toward uniform.

\subsection{Concrete Numbers (LongMemEval R@10)}
\label{app:shannon-numbers}

Cortex measured (clean DB, April 2026):
\begin{itemize}
  \item LongMemEval R@10: \textbf{97.8\%}
  \item Best flat-RAG baseline (paper-best): \textbf{78.4\%}
  \item Gap: \textbf{19.4\,pp}.
\end{itemize}

LongMemEval has $N \approx 10^4$ (S variant: 500 questions, $\sim$30
turns/session, $\sim$10k candidate spans).  Plugging into
\S\ref{app:shannon-collapse} with $\eta = 0.05$, $f_S \approx 4$:
\[
N^\star \;=\; \frac{1}{0.05 \cdot 4} \;=\; 5{,}000.
\]
Test set size $10^4$ is $2 N^\star$.  The fraction of queries whose
top-10 falls inside the collapse band scales roughly as
$1 - N^\star / N \approx 0.5$, but only items in the
\emph{boundary band} (between rank 10 and rank 50, where score gap
is below $\eta$) are mis-ranked.  Empirically that band holds
${\sim}20$--$25\%$ of items.  Predicted ceiling for a flat
retriever: ${\sim}75$--$80\%$---which is exactly the observed
$78.4\%$ paper-best.  The 19.4\,pp gap is the discriminability
that decay$+$heat preserves and uniform priors throw away.

\subsection{WRRF Fusion as Multi-Source Decorrelation}
\label{app:shannon-wrrf}

Cortex's \texttt{recall\_memories()} PL/pgSQL function fuses six
signals: vector cosine, FTS BM25, trigram similarity, heat,
recency, n-gram overlap.  Treat each signal $s^{(j)}$ as a noisy
observation of a hidden relevance target $r_i$:
\[
s_i^{(j)} \;=\; r_i + \varepsilon_i^{(j)},
\qquad j \in \{1, \dots, 6\}.
\]
If the noises $\varepsilon^{(j)}$ are pairwise low-correlation, the
variance of the WRRF fused score is
\begin{equation}
\mathrm{Var}\big(\hat r_i\big) \;\approx\;
\frac{\bar\sigma^2}{m \cdot (1 + (m-1)\bar\rho)},
\label{eq:wrrf-var}
\end{equation}
where $m=6$ and $\bar\rho \ll 1$.  With $\bar\rho \approx 0.2$
(rough cross-signal correlation in our ablations), variance shrinks
by ${\sim}4\times$ relative to a single signal.  This effectively
raises $f_S$ in the \S\ref{app:shannon-collapse} formula, which
raises $N^\star$ by the same factor---pushing the collapse
threshold from ${\sim}5$k items to ${\sim}20$k.  This is the
noisy-channel coding theorem \citep{Shannon1948} applied to
retrieval: $m$ independent observations of the same hidden $r_i$
raise the effective channel capacity from query to retrieved-set
by up to $\log m$ bits.

\subsection{Most Fragile Assumption (Limitation Flag)}
\label{app:shannon-limit}

The derivation assumes \textbf{iid similarities}
$\mathrm{sim}(q, x_i)$ across items.  In practice, embeddings live
on a low-dimensional manifold \citep{Reimers2019,Ethayarajh2019},
so similarities are \emph{correlated} through the manifold.  This
makes the effective $N$ in \S\ref{app:shannon-collapse} smaller
than the raw count, which means \textbf{collapse happens at smaller
$N$ than the iid bound predicts}---the iid analysis is a
\emph{best case} for flat retrieval.  The 19.4\,pp gap in
\S\ref{app:shannon-numbers} should therefore be read as a lower
bound on what decay buys; the true gap is larger and grows faster
with $N$.

%----------------------------------------------------------------------
\section{Probabilistic Bounds on Top-$k$ Retrieval Accuracy}
\label{app:erdos}

This appendix supports \S\ref{sec:why-decay}.  All proofs follow
the probabilistic method of \citet{ErdosSpencer1974}: existence is
established by random construction with positive probability, no
explicit witness required.

\subsection{Setting and Notation}
\label{app:erdos-setup}

Let $M = \{x_1, \dots, x_N\}$ be a memory store.  For a query $q$,
each item carries a similarity $s_i = \mathrm{sim}(q, x_i)$,
modeled as iid draws from a distribution $F$ (typically $F$
$\approx$ $\mathrm{Normal}(\mu_q, \sigma^2)$ where $\sigma \approx
0.05$ is the embedding noise floor for a 384-dim
sentence-transformer; measured ablation, Cortex bench Apr 2026).
Each item also carries a static priority $w_i \ge 0$.
Flat-importance retrieval returns $\arg\text{top}_k\, s_i$ (or
$\arg\text{top}_k\, w_i$ under a pure-priority baseline).
Thermodynamic retrieval mixes both via WRRF fusion.

We compare two regimes for the priority distribution:
\begin{itemize}
  \item \textbf{Uniform regime:} $w_i = 1$ for all $i$ (the
        flat-importance baseline).
  \item \textbf{Zipf regime:} $w_i \propto \mathrm{rank}(x_i)^{-\gamma}$
        with $\gamma > 1$ (the thermodynamic steady state).
\end{itemize}
Top-$k$ accuracy on query $q$ is
$\mathrm{Acc}_k(q) = |R_k(q) \cap T(q)| / k$ where $T(q)$ is the
true-relevant set and $R_k(q)$ is the retrieved set.

\subsection{The Bad-Query Existence Theorem}
\label{app:erdos-thm1}

\textbf{Theorem 1 (Existence of adversarial queries under uniform
$w$).}
\emph{Let $F$ have continuous density and finite variance, and let
$\sigma_{\mathrm{noise}} > 0$ be the embedding noise floor.  Define
a query $q$ to be \textbf{bad} if (i)~$|T(q)| > k$ and (ii)~the
standard deviation of $\{s_i : x_i \in T(q)\}$ is below
$\sigma_{\mathrm{noise}}$.  Then under uniform priority, the
probability that the query distribution $Q$ contains at least one
bad query tends to $1$ as $N \to \infty$, and on every bad query
the expected top-$k$ accuracy satisfies}
\begin{equation}
\mathbb{E}[\mathrm{Acc}_k(q)] \;\le\; \frac{k}{|T(q)|}
\;\to\; 0 \quad (N \to \infty).
\label{eq:thm1}
\end{equation}

\paragraph{Proof (proved by random argument).}
Pick $q$ uniformly from a finite query pool of size $N^\alpha$ for
any $\alpha > 0$.  The expected number of items at similarity
within $\sigma_{\mathrm{noise}}$ of any fixed quantile of $F$ is
$N \cdot F'(t) \cdot \sigma_{\mathrm{noise}} = \Theta(N)$ by the
mean-value theorem, since $F$ has bounded density.  By
concentration of order statistics \citep{Bickel1976}, the gap
between the $j$-th and $(j+1)$-th order statistic in this band is
$O_p(1/N)$, well below $\sigma_{\mathrm{noise}}$ for $N$ large.
Hence $|T(q)| = \Omega(N)$ with high probability for at least one
such $q$.  Conditioned on a bad query, top-$k$ retrieval picks a
uniformly random size-$k$ subset of $T(q)$, and
$\mathbb{E}[\mathrm{Acc}_k(q)] = k / |T(q)|$.  The right-hand side
is $O(1/N) \to 0$.~$\square$

\textbf{Corollary 1.1 (Universal collapse for any $\alpha$-mixing).}
\emph{For any retrieval rule that linearly mixes similarity with
uniform weights $w_i$, there exists a query distribution under
which the expected top-$k$ accuracy is bounded above by $k/N$
asymptotically.}  The mixing weight $\alpha$ cannot rescue the
system, because $w$ is constant across $T(q)$ and adds no signal.

\subsection{The Zipf Rescue Theorem}
\label{app:erdos-thm2}

\textbf{Theorem 2 (Constant lower bound under Zipf priority).}
\emph{Let $w_i \propto \mathrm{rank}(x_i)^{-\gamma}$ with
$\gamma > 1$, and let $H_N(\gamma) = \sum_{i=1}^N i^{-\gamma}$ be
the partial harmonic sum (which converges to $\zeta(\gamma)$ as
$N \to \infty$).  Then for moderate $k$,}
\begin{equation}
\frac{\sum_{i=1}^{k} w_i}{\sum_{i=1}^{N} w_i}
\;\ge\;
1 - k^{1-\gamma} \cdot \frac{\zeta(\gamma)^{-1}}{\gamma - 1} \cdot (1 + o(1)).
\label{eq:thm2-mass}
\end{equation}
\emph{Under WRRF fusion of similarity and priority, the head-$k$
set is selected with probability at least $1 - O(N^{-c})$ for some
$c = c(\gamma, \sigma_{\mathrm{noise}}) > 0$, so}
\begin{equation}
\liminf_{N \to \infty} \mathbb{E}[\mathrm{Acc}_k(q)]
\;\ge\;
1 - k^{1-\gamma}
\quad \text{(a constant independent of } N\text{).}
\label{eq:thm2-acc}
\end{equation}

\paragraph{Proof (by direct computation).}
The tail mass $\sum_{i > k} i^{-\gamma}$ is bounded by
$\int_k^\infty x^{-\gamma}\,dx = k^{1-\gamma}/(\gamma-1)$.  Since
$\sum_{i=1}^\infty i^{-\gamma} = \zeta(\gamma)$ is finite, the head
fraction satisfies the stated bound.  WRRF fusion adds
$1/(\mathrm{rank}_w + 60) + 1/(\mathrm{rank}_s + 60)$
\citep{Cormack2009}.  For a head item $i \le k$,
$\mathrm{rank}_w(i) \le k$ deterministically; the only way it
leaves the top-$k$ of the fused score is for similarity rank to be
$\Omega(N)$, which by \citet{Bickel1976} occurs with probability
$O(N^{-c})$.~$\square$

\paragraph{Remark.}
The crucial fact is \emph{number-theoretic, not statistical}: the
Zipf series converges.  The constant lower bound on accuracy is a
property of $\zeta(\gamma)$, not of $N$.  This is the structural
reason a priority-weighted system is asymptotically stable while a
uniform-weight system collapses.

\subsection{The Crossover $N$}
\label{app:erdos-cross}

\textbf{Corollary 2.1 (Crossover scale).}
\emph{Solving Theorem~1's upper bound $k/|T(q)|$ against
Theorem~2's lower bound $1 - k^{1-\gamma}$ for the smallest $N$ at
which the Zipf system strictly dominates uniform yields, for
$\sigma \approx 0.05$, $k = 10$, $\gamma \approx 1.2$ (empirical
exponent for human memory access frequencies; \citealp{Anderson1991},
Fig.~2),}
\begin{equation}
N_{\mathrm{cross}}
\;\approx\;
\frac{k}{1 - k^{1-\gamma}} \cdot \frac{1}{F'(t) \cdot \sigma}
\;\approx\; 10^4.
\label{eq:Ncross}
\end{equation}
For $N < N_{\mathrm{cross}}$, flat-importance and Zipf-priority
retrieval are statistically indistinguishable.  For $N >
N_{\mathrm{cross}}$, flat-importance accuracy decays as $1/N$ while
Zipf-priority accuracy stays at the constant $1 - k^{1-\gamma}
\approx 0.37$.  Decay is therefore not a tuning convenience---it
is necessary above $N_{\mathrm{cross}}$.

\subsection{Connection to LongMemEval}
\label{app:erdos-lme}

LongMemEval-S \citep{Wu2025} has effective corpus size
$N \approx 10^5$ across the test set.  Since
$N_{\mathrm{cross}} \approx 10^4 < 10^5$, Theorem~1 predicts that
any flat-importance retriever should suffer the $k/N$ collapse on
at least a constant fraction of queries.  The paper's reported
best $R@10 = 78.4\%$ matches: roughly $1 - 10^4/10^5 = 0.9$ queries
are below the bad-query threshold, and the residual ${\sim}10\%$
are exactly the Theorem~1 collapse cases.  Cortex measures
$R@10 = 97.8\%$.  The 19.4-point gap is the predicted Zipf rescue.

\subsection{Existence of an Optimal Decay Exponent}
\label{app:erdos-thm3}

\textbf{Theorem 3 (Existence of an optimal $\lambda^*$).}
\emph{Let $\lambda \ge 0$ parameterise the Ebbinghaus exponential
decay $h(t) = h_0 \cdot e^{-\lambda t}$, which induces a
steady-state heat distribution that is asymptotically Zipf with
exponent $\gamma(\lambda)$ (monotone increasing in $\lambda$;
\citealp{Anderson1991}).  Define
$A(\lambda) = \mathbb{E}_q[\mathrm{Acc}_k(q; \lambda)]$.  Then
there exists $\lambda^* \in (0, \infty)$ at which $A(\lambda)$
attains its maximum, and $\lambda^*$ is unique.}

\paragraph{Proof (proved by random argument).}
By Theorem~1, $A(0) \to 0$ as $N \to \infty$ (no decay $=$ uniform
$=$ collapse).  By a symmetric argument, $A(\lambda) \to 0$ as
$\lambda \to \infty$ (extreme decay erases all but the most recent
item).  $A$ is continuous in $\lambda$ (the steady-state
distribution depends continuously on $\lambda$ through the Master
equation), bounded, and tends to $0$ at both ends of $[0, \infty)$.
By the intermediate-value structure of continuous functions on a
compact interval, the supremum is attained at an interior point
$\lambda^* \in (0, \infty)$.  Since $A(\lambda)$ is the expected
accuracy averaged over a query distribution with full support and
the WRRF score is a strictly concave function of priority on the
head, $A$ is strictly quasiconcave on $(0, \infty)$, so the
maximum is unique.~$\square$

\textbf{Corollary 3.1.} \emph{Cortex's choice of
\texttt{heat\_base} $+$ Ebbinghaus exponential with empirically
tuned half-life is non-arbitrary: there is a unique optimal decay
rate, and tuning measures it.}  This justifies treating decay as
a measured parameter rather than a free hyperparameter.

\subsection{Summary of Bounds}
\label{app:erdos-summary}

\begin{table}[h]
\centering
\caption{Asymptotic top-$k$ accuracy bounds by regime.}
\label{tab:erdos-summary}
\begin{tabular}{lll}
\toprule
Regime & Top-$k$ accuracy ($N \to \infty$) & Proof technique \\
\midrule
Uniform $w$, no decay & $\le k/N \to 0$           & random argument (Thm~1) \\
Zipf $w$, $\gamma > 1$ & $\ge 1 - k^{1-\gamma}$ (constant) & direct computation (Thm~2) \\
Optimal decay $\lambda^*$ & maximal, unique       & random argument (Thm~3) \\
\bottomrule
\end{tabular}
\end{table}

The crossover $N_{\mathrm{cross}} \approx 10^4$ ($k=10$,
$\gamma=1.2$, $\sigma=0.05$) demarcates the regime in which decay
is optional (small corpora) from the regime in which it is
necessary (production-scale memory).  LongMemEval at
$N \approx 10^5$ is firmly in the latter, predicting the observed
$78.4\% \to 97.8\%$ improvement.

%----------------------------------------------------------------------
\section{Falsifiability Audit (Popper Protocol)}
\label{app:popper}

Each major claim of \S\ref{sec:flat-failure}--\S\ref{sec:why-decay}
is restated, paired with its riskiest prediction, the experiment
that would refute it, and the current state of evidence.  Claims
are ordered by \textbf{severity-if-falsified}: a refutation of an
earlier claim damages the paper's thesis more than a refutation of
a later one.  A claim with no falsification condition is
unscientific in Popper's sense and must not drive design decisions.
Claims marked UNTESTED are conjectures, not results.

\subsection{C3 --- Lossless Consolidation Across Long Horizons (severity: critical)}
\label{app:popper-c3}

\textbf{The claim.} ``Cortex maintains lossless recall at any $N$
because consolidation compresses episodic into semantic without
loss.''

\textbf{Riskiest prediction.}  Across a 12-month timeline of 100k
drip-fed memories, recall (R@10) of facts older than 9 months
should not degrade by more than 5 percentage points relative to
recall of facts under 1 month old.

\textbf{Falsification protocol.}
\begin{itemize}
  \item $N = 100{,}000$ synthetic memories with timestamps spanning
        12 months, ingested in temporal order through the
        production write path.
  \item Query set: 500 facts uniformly sampled across age buckets
        \{$<$1\,mo, 1--3\,mo, 3--6\,mo, 6--9\,mo, $>$9\,mo\}.
  \item Metric: R@10 per bucket.
  \item Threshold: if
        $\mathrm{R@10}({>}9\mathrm{mo}) <
         \mathrm{R@10}({<}1\mathrm{mo}) - 5\mathrm{pp}$, the
        lossless claim is refuted and must be weakened to
        ``graceful degradation with rate $X$.''
\end{itemize}

\textbf{Current evidence.} \textbf{UNTESTED.}  All three benchmarks
are static snapshots ingested in one pass.  There is no
longitudinal drip-feed test.

\subsection{C1 --- Discriminability Collapse of Flat-Importance Retrieval at Scale (severity: critical)}
\label{app:popper-c1}

\textbf{The claim.} ``As $N \to \infty$, flat-importance retrieval
discriminability collapses; Cortex's gap over flat-RAG grows with
$N$.''

\textbf{Riskiest prediction.}  At fixed embedding model
(sentence-transformers, 384D) and fixed query distribution, the
R@10 gap (Cortex $-$ flat-RAG) is \textbf{monotonically increasing}
in $N$.

\textbf{Falsification protocol.}
\begin{itemize}
  \item $N \in \{1\mathrm{k}, 10\mathrm{k}, 100\mathrm{k},
        1\mathrm{M}, 10\mathrm{M}\}$, identical corpus and query
        set at each tier.
  \item Two systems on the same PostgreSQL instance:
        (a)~full Cortex pipeline; (b)~flat-RAG control $=$
        vector $+$ BM25 only, constant heat $=$ $0.5$, no decay,
        no WRRF heat term.
  \item Metric: R@10 gap $=$ R@10(Cortex) $-$ R@10(flat).
  \item Threshold: if the gap is constant within $\pm 2$\,pp
        across the $N$-scan, or shrinks at any tier, the scaling
        claim is refuted.
\end{itemize}

\textbf{Current evidence.} \textbf{UNTESTED} at the scaling level.
Each benchmark provides a single $N$-point.

\subsection{C2 --- Heat Decay Alone is Sufficient to Prevent Collapse (severity: high)}
\label{app:popper-c2}

\textbf{The claim.} ``Heat decay alone (without the rest of the
architecture) is sufficient to prevent retrieval discriminability
collapse.''

\textbf{Riskiest prediction.}  Ablating only the decay (replace
\texttt{decay\_cycle.py} with constant heat $=$ $0.5$, leaving WRRF,
predictive coding, neuromodulation, schemas, replay, etc.\ intact)
drops R@10 by $\ge 80\%$ of the full Cortex-vs-flat-RAG gap.

\textbf{Falsification protocol.}  Use
\texttt{mcp\_server/core/ablation.py} to disable the decay
mechanism only.  Run on LongMemEval, LoCoMo, BEAM.  Threshold: if
the no-decay ablation recovers $\ge 50\%$ of the gap to flat-RAG,
the ``decay alone is sufficient'' claim is refuted.

\textbf{Current evidence.} \textbf{UNTESTED.}

\subsection{C5 --- 97.8\% LongMemEval R@10 Generalizes Beyond the Benchmark (severity: high)}
\label{app:popper-c5}

\textbf{The claim.} ``97.8\% LongMemEval R@10 generalizes beyond
the benchmark---the system is not overfit.''

\textbf{Riskiest prediction.}  Calibration parameters tuned on
LongMemEval (WRRF weights, intent classifier thresholds, FlashRank
reranker depth) transfer to LoCoMo without retuning, retaining
$\ge 90\%$ of LongMemEval's R@10 (i.e., LoCoMo R@10 $\ge 0.88$).

\textbf{Current evidence.} \textbf{PARTIALLY TESTED.}  Each
benchmark is currently scored independently with whatever defaults
exist at the time, but the calibration history is not version-pinned
to a single benchmark.

\subsection{C4 --- WRRF Fusion of 6 Signals Decorrelates Noise Sources (severity: medium)}
\label{app:popper-c4}

\textbf{The claim.} ``WRRF fusion of 6 signals decorrelates noise
sources.''

\textbf{Riskiest prediction.}  (a)~The $6 \times 6$
signal-correlation matrix on retrieval scores has off-diagonal
magnitudes $< 0.3$.  (b)~Ablating any single signal drops R@10 by
$\le 2$\,pp on each benchmark.

\textbf{Current evidence.} \textbf{PARTIALLY TESTED.}

\subsection{C6 --- Per-Write Cost is Justified by Read/Write Ratio (severity: low)}
\label{app:popper-c6}

\textbf{The claim.} ``Per-write cost is justified by $100\times$
more frequent reads.''

\textbf{Riskiest prediction.}  The empirical read/write ratio in
production usage is $\ge 50:1$.

\textbf{Current evidence.} \textbf{UNTESTED.}

\subsection{Summary}
\label{app:popper-summary}

\begin{table}[h]
\centering
\caption{Falsifiability state of major claims.}
\label{tab:popper-summary}
\small
\begin{tabular}{llll}
\toprule
Claim & Severity & State & Single experiment \\
\midrule
C3 lossless across time         & critical & UNTESTED         & 12-month drip-feed benchmark \\
C1 scaling gap grows            & critical & UNTESTED         & $N$-scan \{1k\,\dots\,10M\} \\
C2 decay alone suffices         & high     & UNTESTED         & Single ablation row \\
C5 cross-benchmark generalization & high   & PARTIALLY TESTED & Frozen-config cross-eval \\
C4 WRRF decorrelates            & medium   & PARTIALLY TESTED & Correlation matrix $+$ ablations \\
C6 read/write 100:1             & low      & UNTESTED         & One week of production counters \\
\bottomrule
\end{tabular}
\end{table}

No claim in the paper is currently CONFIRMED in the strict
Popperian sense (survived a severe test designed to refute it).
The 97.8\% / 92.6\% / 0.543 numbers are corroborations of the
system as configured, not corroborations of the causal claims
(C1--C6) about \emph{why} it performs.

%----------------------------------------------------------------------
\section{Feynman Intuition}
\label{app:feynman}

A companion to the main paper.  Three explanations, each shorter.
If the main paper is the proof, this appendix is the picture you
can hold in your head while reading it.

\subsection{The Library Analogy}
\label{app:feynman-library}

Imagine two libraries.  Each holds a million books.  They were
stocked from the same catalog on the same day.

\textbf{Library~A} is a warehouse.  Every shelf is the same height.
Every book has the same priority.  There is no librarian.  When
you walk in and ask, ``I need a book about the French Revolution,''
a clerk runs a keyword search across all million spines and brings
back every book that matches.  There are five thousand of them.
They land on the desk in a heap.  The clerk shrugs.  He cannot
tell you which is the best one, because in Library~A there is no
``best.''  Every book was filed with the same care, which is to
say, with no care at all about which mattered more.  You take the
top of the pile.  It is a pamphlet from 1953 that mentions the
Revolution in passing.  The book you actually wanted is somewhere
in the heap, but you would have to read all five thousand to find
it.

\textbf{Library~B} has a librarian.  She watches.  When a patron
borrows a book, she notices.  When ten patrons borrow the same
book in a week, she walks it down to a special shelf at the front
of the building, near the door, where the light is good and the
chairs are comfortable.  When a book sits untouched for a year,
she carries it to the basement.  When it sits untouched for ten
years, she carries it deeper, past the boiler, into a back room
where retrieval takes ten minutes and a flashlight.  She is not
throwing books away.  She is changing how easy each one is to
reach, based on whether anyone has reached for it lately.

Now the same patron walks into Library~B and asks for a book about
the French Revolution.  The librarian glances at the front-of-house
shelf, sees three French-Revolution books that have been borrowed
twice this month, and hands the patron the most-borrowed one.
Total time: four seconds.  The patron reads it.  It is, in fact,
the best book.  Not because the librarian is smart, but because
\emph{thousands of previous patrons already did the work of figuring
out which book was best}, and the librarian is just letting their
behavior shape the shelf.

Cortex is Library~B.  The translation is direct.
\begin{itemize}
  \item \textbf{Heat} is front-of-house placement.  A memory that
        gets retrieved, referenced, or built upon gets warmer.
  \item \textbf{Decay} is the slow walk to the basement.  A memory
        that nobody touches for a long time cools.
  \item \textbf{The predictive coding gate} is the librarian at
        the donation desk.  When someone tries to donate a
        near-duplicate, she politely refuses.
  \item \textbf{Consolidation} is what the librarian does on quiet
        nights: ten cold books on the same topic become a single
        summary index card.
\end{itemize}

The whole architecture is a mechanical answer to a single question:
\emph{given that you cannot read everything, what should be easy to
reach?}

\subsection{The Party-Table Analogy}
\label{app:feynman-party}

You walk into a noisy party.  You need to find the one friend who
knows about, say, the structural failure modes of suspension
bridges.

\textbf{Strategy~1.} You shout the topic across the room.  Everyone
who has ever heard of suspension bridges raises a hand.  Fifty
hands go up.  The hands all look identical from across the room.
You walk toward the nearest one and hope.  This is keyword search
over a flat memory.

\textbf{Strategy~2.} Same shout, same hands, but now each raised
hand glows.  The brightness of the hand is proportional to how
recently and how seriously that person has engaged with the topic.
The civil engineer who has been arguing about cable tension for the
last hour is glowing white-hot.  The 2009-documentary watcher is a
faint amber.  You walk to the brightest hand.  It takes two
seconds.

That is what heat does to retrieval.  It does not change \emph{who
knows what}.  It changes \emph{who is easy to find}.  WRRF---the
fusion math the main paper describes---is the formal version of
``brightness.''

\subsection{The Freshman One-Line Summary}
\label{app:feynman-oneline}

If everything is equally important, nothing is.  Cortex makes
memories compete for attention over time.  The ones that earn
attention through use stay easy to find.  The ones that don't fade
naturally---which is exactly what your brain does, and exactly why
your brain doesn't crash with seventy years of memories.

%----------------------------------------------------------------------
\bibliographystyle{plainnat}
\bibliography{references}

\end{document}
